\scanpage{44}
\uline{Question} Si $(S,Z)$ est un système $n$-régulier de sommets dans $E$, $Z$ est-il déterminé de façon unique par la donnée de la seule partie $S$ de $E$ ? On trouvera ceci : si $(S,\omega)$ est un $n$-polygône régulier dans $E$, le stabilisateur \inserted{$G$} de $S$ dans $\mathrm{Aff}(E)$ (je dis bien de \inserted{la \uline{partie}} $S$ de $E$, non du couple $(S,\omega)$) est égal au groupe $\operatorname{Aut}(E,S,\omega)$, sauf dans les cas $\alpha=2$, \deleted{où on trouve un groupe isom} \inserted{ou $\alpha=-2$ ; dans ces cas, identifiant} $G$ à un sous-groupe de $\unsure{\mathfrak S_S|\mathcal S_S|G_S}$, on trouve $\mathrm{Aff}(S,Z)$.
% slanted
\margin{\uline{à vérifier} !}

Donc la question devient purement combinatoire, savoir : une structure quasi polygonale comb. \inserted{$Z$} sur $S$ est-elle déterminée par la seule connaissance de $\mathrm{Aff}^{\pm}(S)\subset\mathfrak S_S$, ou par celle de $\mathrm{Aff}(S)\subset\mathfrak S_S$ ? Dans le premier cas, $Z$ est déterminé comme le plus grand ss-groupe cyclique \uline{distingué} $\neq G=\mathrm{Aff}^{\pm}(S)$. [car si $\sigma\in G\smallsetminus Z$, on voit facilem$^{\rm t}$ que
% archive sheet number: 43
